% CCSS 7.SP.B.3 — Comparing Two Populations Visually (FL/IN/MN/VA: multiple display emphasis)
\section{Comparing Two Populations Visually}

\begin{learningGoals}
\begin{itemize}
    \item Compare two data sets using multiple display types
    \item Draw informal conclusions about differences between populations
\end{itemize}
\end{learningGoals}

\begin{conceptBox}{Tools for Visual Comparison}
Use these displays to compare two populations side by side:
\begin{itemize}[leftmargin=*]
    \item \textbf{Dot plots} — show every data point; great for small data sets.
    \item \textbf{Histograms} — group data into intervals; show the shape of the distribution.
    \item \textbf{Box plots} — show median, quartiles, and range at a glance.
    \item \textbf{Stem-and-leaf plots} — keep exact values while showing shape.
    \item \textbf{Circle graphs} — compare proportions (parts of a whole) between groups.
    \item \textbf{Frequency tables} — organize counts by category or interval.
\end{itemize}
Look at both the \textbf{center} (mean/median) and the \textbf{spread} (range, IQR) when comparing.
\end{conceptBox}

\begin{workedExample}{Comparing with a Back-to-Back Stem-and-Leaf Plot}
Two classes recorded test scores:

\begin{center}
\begin{tabular}{r|c|l}
\textbf{Class A} & \textbf{Stem} & \textbf{Class B} \\
\hline
$8\;5$ & $6$ & $2\;7$ \\
$9\;6\;3\;0$ & $7$ & $1\;4\;8$ \\
$7\;5\;2$ & $8$ & $0\;3\;5\;9$ \\
$4$ & $9$ & $2\;6$
\end{tabular}
\end{center}

\textbf{Class A:} Median $\approx 76.5$, range $= 94 - 65 = 29$.

\textbf{Class B:} Median $\approx 80$, range $= 96 - 62 = 34$.

\answerHighlight{Class B has a higher center, but Class A's scores are more tightly clustered.}
\end{workedExample}

\mascotSays{One display tells part of the story. Try two or three different displays to see the full picture!}

\begin{practiceBox}{Comparing Populations Visually}
\resetProblems

\prob Two teams' points: Team~X has a mean of $80$ and range of $15$. Team~Y has a mean of $85$ and range of $30$. Compare them.
\answerExplain{Team Y averages higher but is more spread out}{Team Y's higher mean shows better scoring on average, but its larger range means more inconsistency game to game.}

\prob When is a box plot better than a dot plot for comparing two groups?
\answerExplain{When data sets are large}{Box plots summarize key statistics (median, quartiles, range) without showing every point, making them easier to read when there are many data values.}

\prob A circle graph shows Class~A spends $40\%$ of free time on sports and Class~B spends $25\%$. What can you conclude?
\answerExplain{Class A devotes a larger share to sports}{Circle graphs compare proportions. A larger sector means a bigger fraction of the whole. Class A spends more of their free time on sports.}

\prob Two groups have overlapping box plots. Does this mean the groups are the same?
\answerExplain{Not necessarily}{Overlap means some members are similar, but the centers may still differ. Compare medians and IQRs to see if the difference is meaningful.}

\prob Name a situation where a frequency table is more useful than a histogram.
\answerExplain{Categorical data (e.g., favorite color)}{Histograms need numerical intervals. For categories like favorite color, a frequency table counts each category directly.}

\end{practiceBox}
